% TEMPLATE-MILC-v3.tex - plantilla maestra UNICA de las guias MILC v3.
%
% La usan las tres familias de guias, cada una con su driver:
%   · Grados 6-11  -> scripts/build-guias-g11.py
%   · Bebras       -> scripts/build-guias-bebras.py
%   · Semillero    -> scripts/build-guias-semillero.py
%
% Antes habia tres copias casi identicas (~400 lineas iguales, 6 distintas):
% cada mejora habia que aplicarla tres veces, y en la practica no se hacia
% (el motor de recursos vivia solo en dos de ellas). Lo que varia entre
% familias esta parametrizado en 6 placeholders que cada driver llena
% (se nombran aqui SIN su sintaxis real: el builder reemplaza por texto plano,
% tambien dentro de los comentarios, y un valor multilinea rompe el preambulo):
%
%   HEADER_IZQ        encabezado izquierdo de pagina
%   HEADER_DER        encabezado derecho de pagina
%   PORTADA_ETIQUETA  rotulo sobre el numero grande (GRADO/MOMENTO/MODULO)
%   PORTADA_SERIE     linea de serie de la portada
%   PORTADA_ID        identificador de la guia en la portada
%   PORTADA_CIRCULO   el circulito de grado (°), o vacio si no aplica
%
% El resto de claves las documenta content/guias/_SCHEMA.md. El script valida
% que no queden placeholders sin reemplazar antes de escribir el archivo final.

\documentclass[11pt,letterpaper]{article}
\usepackage[margin=2.0cm]{geometry}
\usepackage[table]{xcolor}
\usepackage{fontspec}
\usepackage[spanish]{babel}
\usepackage{tikz}
% Librerías TikZ para los diagramas de las guías (bloque `recursos:`):
%  · babel  → babel en español vuelve activos caracteres como «>», y TikZ los usa
%             en sus opciones (>=stealth, ->). Sin esta librería, cualquier
%             diagrama con flechas revienta con «\language@active@arg> has an
%             extra }». Debe ir ANTES de que se dibuje nada.
%  · positioning        → sintaxis «below left=2pt and 3pt».
%  · decorations.pathreplacing → llaves ({brace}) para señalar rangos.
\usetikzlibrary{babel,positioning,decorations.pathreplacing}
\usepackage{graphicx}
% Circuitos: el trabajo de electrónica se hace hoy con micro:bit y sus sensores
% integrados (MakeCode, sin cableado), así que no se necesitan esquemáticos.
% Si algún día entran montajes con protoboard, basta descomentar esta línea:
% los diagramas .tex/.tikz declarados en `recursos:` ya se incrustan solos.
% \usepackage{circuitikz}
\usepackage[most]{tcolorbox}
\usepackage{tabularx}
\usepackage{array}
\usepackage{fancyhdr}
\usepackage{titlesec}
\usepackage[document]{ragged2e}  % bandera a la derecha en todo el cuerpo
\usepackage{enumitem}
\usepackage{microtype}

% Helvetica Neue no trae la flecha U+2192, asi que cada «→» escrito literalmente
% en el YAML se compilaba a NADA, en silencio (462 flechas en 61 guias: rutas del
% tipo «paleta → tipografia → lockup» salian rotas). Mapeamos el caracter a su
% version matematica, que si tiene glifo. Asi el autor puede seguir escribiendo
% «→» comodamente en el YAML.
\usepackage{newunicodechar}
\newunicodechar{→}{\ensuremath{\rightarrow}}
% Mismo problema con los simbolos de marca y vinetas: el «✓» de «una palomita ✓»
% salia como cuadrito vacio en el PDF de todas las guias que lo usaban.
\usepackage{pifont}
\newunicodechar{✓}{\ding{51}}
\newunicodechar{✗}{\ding{55}}
\newunicodechar{✔}{\ding{52}}
\newunicodechar{•}{\textbullet}
\newunicodechar{≥}{\ensuremath{\geq}}
\newunicodechar{≤}{\ensuremath{\leq}}

\setmainfont{Helvetica Neue}
\setsansfont{Helvetica Neue}
\setlength{\parindent}{0pt}
\setlength{\parskip}{6pt}
% Sin particion de palabras: en 6.º-7.º una palabra cortada por guion al final
% de linea («Comuni-cacion») frena la lectura mucho mas que una linea corta.
\hyphenpenalty=10000
\exhyphenpenalty=10000
\pretolerance=2000
\tolerance=3000
\setlength{\emergencystretch}{3em}
\linespread{1.10}
\renewcommand{\arraystretch}{1.25}
\setlist{leftmargin=5mm,itemsep=1mm,topsep=1mm}
\newcolumntype{Y}{>{\RaggedRight\arraybackslash}X}

\definecolor{milcMagenta}{HTML}{E6007E}
\definecolor{milcVino}{HTML}{5A0038}
\definecolor{milcTurquesa}{HTML}{0093A5}
\definecolor{milcVerde}{HTML}{58A923}
\definecolor{milcMostaza}{HTML}{E5B400}
\definecolor{milcOcre}{HTML}{B86600}
\definecolor{milcNegro}{HTML}{241816}
\definecolor{milcGris}{HTML}{F7F7F4}
\definecolor{milcLinea}{HTML}{E8E2DD}
\definecolor{milcGradePrimary}{HTML}{EA580C}
\definecolor{milcGradeDark}{HTML}{9A3412}
\definecolor{milcGradeSoft}{HTML}{FFEDD5}
\definecolor{milcGradeAccent}{HTML}{CFE7FF}
\definecolor{milcGradeText}{HTML}{FFFFFF}
\color{milcNegro}

\pagestyle{fancy}
\fancyhf{}
\lhead{\small Semillero de Investigación · Robótica y sistemas embebidos}
\rhead{\small Módulo 4 · Evaluación liberadora · MILC v3}
\cfoot{\small\thepage}
\renewcommand{\headrulewidth}{0pt}

\newtcolorbox{titlebox}[1]{
  enhanced,colback=#1,colframe=#1,arc=7mm,boxrule=0pt,
  left=7mm,right=7mm,top=3mm,bottom=3mm,
  before skip=8pt,after skip=8pt,
  fontupper=\sffamily\bfseries\Large\color{white}
}

\newtcolorbox{softbox}[3]{
  enhanced,colback=#2,colframe=#1,borderline west={4pt}{0pt}{#1},
  arc=2mm,boxrule=.4pt,left=4mm,right=4mm,top=3mm,bottom=3mm,
  before skip=6pt,after skip=8pt,title={#3},coltitle=white,
  colbacktitle=#1,fonttitle=\sffamily\bfseries,fontupper=\RaggedRight,
  attach boxed title to top left={xshift=3mm,yshift=-2mm},
  boxed title style={arc=2mm, boxrule=0pt}
}

\newtcolorbox{drawbox}[2]{
  enhanced,colback=white,colframe=#1,arc=4mm,boxrule=1pt,
  left=5mm,right=5mm,top=4mm,bottom=4mm,before skip=8pt,after skip=8pt,
  title={#2},coltitle=white,colbacktitle=#1,fonttitle=\sffamily\bfseries,
  fontupper=\RaggedRight,
  attach boxed title to top left={xshift=4mm,yshift=-2mm},
  boxed title style={arc=3mm, boxrule=0pt}
}

\newtcolorbox{infoband}[1]{
  enhanced,colback=#1!8,colframe=#1!50,arc=2mm,boxrule=.5pt,
  left=4mm,right=4mm,top=2mm,bottom=2mm,before skip=4pt,after skip=4pt,
  fontupper=\small\RaggedRight
}

% Puente narrativo entre secciones (contrato editorial Conéctate).
% Da continuidad: "Acabas de X. Ahora vas a Y porque Z."
% Visualmente: banda izquierda en color del grado, texto en itálica pequeña.
\newtcolorbox{puentebox}{
  enhanced,colback=milcGradeSoft,colframe=milcGradeDark,
  borderline west={3pt}{0pt}{milcGradeDark},
  arc=0mm,boxrule=0pt,left=5mm,right=4mm,top=2.5mm,bottom=2.5mm,
  before skip=4pt,after skip=8pt,
  fontupper=\itshape\small\color{milcNegro}\RaggedRight
}
% La flecha va en modo matematico a proposito: Helvetica Neue NO tiene el glifo
% U+2192, asi que \textrightarrow se compilaba a NADA («Missing character: There
% is no → in font Helvetica Neue») y el puente salia sin flecha. En modo
% matematico la flecha viene de la fuente matematica y si se dibuja.
\newcommand{\puente}[1]{%
  \begin{puentebox}%
    \textcolor{milcGradeDark}{$\rightarrow$}\hspace{2mm}#1%
  \end{puentebox}%
}

\newcommand{\linead}[1]{\noindent\color{milcLinea}\rule{#1}{0.5pt}\color{milcNegro}}
\newcommand{\respuesta}[1]{\par\vspace{2mm}\foreach \n in {1,...,#1}{\noindent\color{milcLinea}\rule{\linewidth}{0.5pt}\par\vspace{5mm}}\color{milcNegro}}
\newcommand{\checkbox}{\textcolor{milcGradeDark}{\fbox{\phantom{x}}}\hspace{5pt}}

\newcommand{\phasecard}[4]{
\begin{tcolorbox}[enhanced,colback=#3,colframe=#2,arc=3mm,boxrule=.5pt,height=3.05cm,valign=center,halign=center,left=1.5mm,right=1.5mm,top=2mm,bottom=2mm]
{\sffamily\bfseries\fontsize{22}{24}\selectfont\textcolor{#2}{#1}}\par
{\sffamily\bfseries\small #4}
\end{tcolorbox}
}

% ── Piezas del rediseno 2026 (legibilidad 6.º-11.º) ───────────────────
% Banda de estacion: reemplaza el titlebox macizo. Titulo a la izquierda,
% dato practico (tiempo/modalidad) a la derecha, en una sola linea.
\newcommand{\milcestacion}[2]{%
\begin{tcolorbox}[enhanced,colback=#1,colframe=#1,arc=2mm,boxrule=0pt,
  left=5mm,right=5mm,top=2.6mm,bottom=2.6mm,before skip=11pt,after skip=7pt]
{\sffamily\bfseries\fontsize{15}{18}\selectfont\color{white}#2}
\end{tcolorbox}}

% Tarjeta de paso numerada: sustituye las tablas «Pilar / Aplicacion», que
% eran formato de consulta docente, no de estudiante. El numero va en un
% circulo sobre el borde para que la secuencia se lea de un vistazo.
\newcommand{\milcpaso}[3]{%
\begin{tcolorbox}[enhanced,colback=white,colframe=milcLinea,arc=1.5mm,boxrule=.9pt,
  left=14mm,right=4mm,top=3mm,bottom=3mm,before skip=4pt,after skip=4pt,
  fontupper=\RaggedRight,
  overlay={\node[anchor=north west,circle,fill=#2,text=white,inner sep=0pt,
    minimum size=7.4mm,font=\sffamily\bfseries\small]
    at ([xshift=3.6mm,yshift=-3.2mm]frame.north west) {#1};}]
#3
\end{tcolorbox}}

% Casilla marcable de tamano util con lapiz (la anterior era \fbox{\phantom{x}},
% de alto variable segun el texto que la seguia).
\newcommand{\milcmarca}{\framebox[3.8mm]{\rule{0pt}{3.4mm}}}

% Fila marcable de lista suelta (checks de la estacion 1).
\newcommand{\milccheckrow}[1]{%
\par\vspace{1.8mm}\noindent
\makebox[6.4mm][l]{\raisebox{-.55mm}{\textcolor{milcNegro}{\milcmarca}}}%
\parbox[t]{\dimexpr\linewidth-6.4mm\relax}{\RaggedRight #1}\par}

% Celda marcable dentro de la rubrica (tabla de autoevaluacion). Misma
% sangria francesa, pero pensada para vivir dentro de una columna X.
\newcommand{\milcopcion}[1]{%
\makebox[5.2mm][l]{\raisebox{-.55mm}{\textcolor{milcNegro}{\milcmarca}}}%
\parbox[t]{\dimexpr\linewidth-5.2mm\relax}{\RaggedRight #1}}

% ── Recursos visuales de la guía ──────────────────────────────────────
% Declarados en el bloque `recursos:` del YAML y emitidos por el builder.
% \guiaFigura   → imagen rasterizada o PDF (foto, carta celeste, montaje).
% \guiaDiagrama → fuente TikZ/circuitikz (.tex) incrustada como vector:
%                 es la vía para circuitos y esquemáticos editables.
% #1 = ruta relativa al .tex · #2 = pie (puede ir vacío)
\newcommand{\guiaFigura}[2]{%
\begin{center}
\includegraphics[width=0.88\linewidth,height=0.38\textheight,keepaspectratio]{#1}\\[1.5mm]
{\footnotesize\itshape\textcolor{milcNegro!75}{#2}}
\end{center}
}

\newcommand{\guiaDiagrama}[2]{%
\begin{center}
\input{#1}\\[1.5mm]
{\footnotesize\itshape\textcolor{milcNegro!75}{#2}}
\end{center}
}

\begin{document}
\thispagestyle{empty}

% ── Portada ───────────────────────────────────────────────────────────
% Antes: un rectangulo de color a pagina completa. Se veia bien en pantalla,
% pero en la fotocopiadora del colegio se comia el toner y salia gris sucio.
% Ahora la banda ocupa el tercio superior y el resto va en blanco.
\begin{tikzpicture}[remember picture,overlay]
  \path[fill=milcGradePrimary]
    (current page.north west) rectangle ([yshift=-10.6cm]current page.north east);

  \node[anchor=north west,text=milcGradeText] at ([xshift=2.0cm,yshift=-1.55cm]current page.north west)
    {\sffamily\bfseries\fontsize{12}{14}\selectfont MÓDULO};
  \node[anchor=north west,text=milcGradeText,opacity=.85] at ([xshift=2.0cm,yshift=-2.35cm]current page.north west)
    {\sffamily\bfseries\fontsize{8.5}{10}\selectfont SEMILLERO DE INVESTIGACIÓN · CONECTATE};
  \node[anchor=north east,text=milcGradeText,opacity=.85,align=right] at ([xshift=-2.0cm,yshift=-1.55cm]current page.north east)
    {\sffamily\bfseries\fontsize{9}{12}\selectfont I.E. Sor María Juliana\\Cartago, Valle del Cauca};

  \node[anchor=west,text=milcGradeText] (gradonum) at ([xshift=1.85cm,yshift=-7.1cm]current page.north west)
    {\sffamily\bfseries\fontsize{112}{112}\selectfont 4};
  % El circulito de grado (°) solo aplica a las guias de grado («11°»); en
  % Bebras y Semillero el numero grande es un momento/modulo, no un grado.
  % Se posiciona relativo al nodo `gradonum`, no en coordenadas absolutas.
  

  \node[anchor=west,text=milcGradeText,text width=11.4cm,align=left] at ([xshift=1.5cm,yshift=0.5cm]gradonum.east)
    {
     {\sffamily\bfseries\fontsize{9}{11}\selectfont\MakeUppercase{Módulo 4 · Fase MILC: Evaluación liberadora · 3 h de clase}}\\[3mm]
     {\sffamily\bfseries\fontsize{21}{25}\selectfont\setlength{\baselineskip}{27pt}Probar, medir, mejorar ---\\[4pt]de la espiral del\\[4pt]tiempo Misak\\[4pt]al prototipo que evoluciona}\\[3.5mm]
     {\sffamily\bfseries\fontsize{15}{18}\selectfont Robótica y sistemas embebidos}
    };
\end{tikzpicture}

\vspace*{9.4cm}
\vfill

{\sffamily\bfseries\fontsize{8}{10}\selectfont\textcolor{milcGradeDark}{LO QUE TE LLEVAS HOY}}

\begin{tcolorbox}[enhanced,colback=white,colframe=milcNegro,arc=2mm,boxrule=1.6pt,
  left=5mm,right=5mm,top=4mm,bottom=4mm,before skip=3pt,after skip=12pt,
  fontupper=\RaggedRight\fontsize{11}{15.5}\selectfont]
Tu \textbf{prototipo evolucionado, con su historia documentada}: el \textbf{registro de pruebas, fallas y depuración} (qué falló, por qué, qué cambiaste); una \textbf{versión mejorada} con una \textbf{iteración sostenida en evidencia} (mediste antes y después); y una \textbf{nota honesta de límites}: qué NO hace tu aparato, y \textbf{a quién le sirve y a quién no}. Con esto cierras tu segundo proyecto de investigación completo ---y abres la pregunta del siguiente.
\end{tcolorbox}

\begin{tcolorbox}[enhanced,colback=milcGris,colframe=milcGris,arc=2mm,boxrule=0pt,
  left=5mm,right=5mm,top=3.5mm,bottom=3.5mm,after skip=12pt]
{\sffamily\bfseries\fontsize{8}{10}\selectfont\textcolor{milcNegro!55}{ESTA GUÍA ES DE}}\\[3mm]
\begin{tabularx}{\linewidth}{X p{3.2cm} p{3.2cm}}
\linead{\linewidth} & \linead{3.0cm} & \linead{3.0cm}\\[-1mm]
{\fontsize{8.5}{10}\selectfont\textcolor{milcNegro!55}{Nombre completo}} &
{\fontsize{8.5}{10}\selectfont\textcolor{milcNegro!55}{Grupo}} &
{\fontsize{8.5}{10}\selectfont\textcolor{milcNegro!55}{Fecha}}\\
\end{tabularx}
\end{tcolorbox}

\vfill

\noindent\textcolor{milcLinea}{\rule{\linewidth}{1pt}}\\[2mm]
{\sffamily\bfseries\fontsize{9.5}{12}\selectfont Autor: Dr. Álvaro Cárdenas Orozco}\\[1mm]
{\sffamily\fontsize{8.5}{11}\selectfont\textcolor{milcNegro!55}{Metodología MILC · apertura ancestral · pensamiento computacional · triángulo Dussel–estoicismo–Floridi}}

\newpage

\section*{Apertura: Probar, medir, mejorar --- de la espiral del tiempo Misak al prototipo que evoluciona}

\begin{softbox}{milcGradeDark}{milcGradeSoft}{Saber ancestral}
\textbf{Para el pueblo Misak ---del Cauca, vecino de nuestro Valle--- el tiempo no es una línea recta que va y no vuelve.} El tiempo es una \textbf{espiral}: un ir y venir que se enrolla y se desenrolla, y que siempre puede \textbf{regresar al centro}, al origen (Misak-Colombia; Artesanías de Colombia). Por eso los Misak no le dan la espalda a lo vivido: \textbf{para caminar hacia adelante, miran hacia atrás} ---entienden lo que pasó, escuchan a los mayores, y con eso avanzan. Su símbolo es el \textbf{Kuarimpele}, una cinta tejida en espiral desde un centro. Fíjate en la sabiduría: no se progresa olvidando y empezando de cero cada vez; se progresa \textbf{volviendo a mirar lo que se hizo, para hacerlo mejor en la próxima vuelta}. Eso es, exactamente, \textbf{evaluar}. Un buen investigador no tira su prototipo y arranca otro desde cero: lo mira de frente ---qué falló, qué aprendió, qué mejoraría--- y con eso da la \textbf{siguiente vuelta de la espiral}, un poco más arriba. Mejorar no es rehacer: es volver con lo aprendido. \textbf{El Misak avanza en espiral; tú también, cada vez que evalúas antes de seguir.}
\end{softbox}

\begin{softbox}{milcTurquesa}{milcGris}{Saber contemporáneo}
Un aparato no termina cuando «funciona»: termina cuando lo has \textbf{evaluado con honestidad} y lo has hecho \textbf{mejor}. La fase final tiene cuatro oficios. \textbf{(1) Depurar} no es adivinar: es \emph{hallar por qué falla}. Se lee el síntoma (\emph{¿qué hizo mal, en qué escenario?}), se piensa \textbf{una} causa probable, se prueba, y si no era, se piensa otra ---una a la vez, como el que busca la gotera siguiendo el agua. \textbf{(2) El ciclo de mejora:} probar → observar dónde falla → \textbf{cambiar una sola cosa} → volver a probar. Cambiar dos a la vez y no sabrás cuál la arregló. \textbf{(3) Evaluar el desempeño con evidencia:} define un criterio (\emph{¿cuántos de los 5 escenarios pasa?}), mídelo \textbf{antes} y \textbf{después} de tu mejora, y muestra el avance con números ---\emph{``pasaba 3 de 5, ahora pasa 5 de 5''}---, no con «quedó mejor». \textbf{(4) Límites y ética:} ningún aparato lo hace todo. Reconoce qué \textbf{no} puede (el sensor es flojo, probaste pocos casos) y ---clave--- \textbf{a quién sirve y a quién no}: una alarma ayuda, pero no reemplaza el cuidado; venderla como garantía sería engañar. \emph{Un aparato honesto dice lo que hace y lo que no.}
\end{softbox}

\begin{softbox}{milcMostaza}{milcGris}{Pregunta-puente}
\textit{El Misak no avanza olvidando: da la vuelta de la espiral mirando lo que ya caminó; \textbf{¿por qué mejorar tu prototipo ---mirar sus fallas de frente e iterar--- te lleva más lejos que rehacerlo de cero}, y por qué reconocer sus límites lo hace más confiable, no menos?}
\end{softbox}

\begin{softbox}{milcVerde}{milcGris}{Saber y saber hacer}
Al terminar la sesión: (1) podrás \textbf{identificar} las fallas de tu prototipo y sus causas probables (depurar); (2) sabrás \textbf{explicar} el ciclo de mejora, cómo se evalúa el desempeño con evidencia y qué exige la ética del aparato; (3) podrás \textbf{evaluar} tu prototipo y \textbf{crear} una versión mejorada con una iteración documentada; (4) habrás \textbf{reconocido} los límites de tu aparato y a quién sirve y a quién no.
\end{softbox}



\small
\begin{tabularx}{\textwidth}{>{\bfseries}p{3.0cm}Y}
\rowcolor{milcGradePrimary!12}Autor & Dr. Álvaro Cárdenas Orozco\\
DBA & Formula preguntas investigables, produce evidencia con método y comunica hallazgos reconociendo sus límites (investigación estudiantil STEM, transversal 6°--11°)\\
Referentes & micro:bit · MakeCode · Continuidad con la unidad de 8° P2 (guías 8-2-1 a 8-2-10)\\
Producto final & Tu \textbf{prototipo evolucionado, con su historia documentada}: el \textbf{registro de pruebas, fallas y depuración} (qué falló, por qué, qué cambiaste); una \textbf{versión mejorada} con una \textbf{iteración sostenida en evidencia} (mediste antes y después); y una \textbf{nota honesta de límites}: qué NO hace tu aparato, y \textbf{a quién le sirve y a quién no}. Con esto cierras tu segundo proyecto de investigación completo ---y abres la pregunta del siguiente.\\
\end{tabularx}
\normalsize

\puente{Escuchaste (módulo 1), diseñaste (módulo 2) y construiste (módulo 3). Hoy \textbf{cierras la espiral}: evalúas tu prototipo de frente, lo haces mejor con evidencia, y reconoces con honestidad qué hace y qué no.}

\milcestacion{milcGradeDark}{Ruta MILC: así vamos a aprender}
\begin{tabularx}{\textwidth}{>{\centering\arraybackslash}X>{\centering\arraybackslash}X>{\centering\arraybackslash}X>{\centering\arraybackslash}X}
\phasecard{1}{milcVerde}{milcGris}{Escucha\\[-1mm]\small Observo, escucho y pregunto.} &
\phasecard{2}{milcTurquesa}{milcGris}{Entender\\[-1mm]\small Sistematizo y ordeno.} &
\phasecard{3}{milcMagenta}{milcGris}{Hacer\\[-1mm]\small Construyo y aplico.} &
\phasecard{4}{milcVino}{milcGris}{Cierre\\[-1mm]\small Evalúo y reflexiono.}
\end{tabularx}


\puente{Antes de mejorar, mira atrás como el Misak: saca tu prototipo del módulo 3 y sus 5 pruebas. Ahí está lo que hay que entender para avanzar.}

\milcestacion{milcVerde}{Estación 1 de 3 · Escucha}

\begin{softbox}{milcVerde}{milcGris}{Escena para escuchar}
\textbf{Actividad 1 · IDENTIFICA --- Mirar atrás las fallas} (40 min · individual + grupo). \textbf{Momento A (10 min) --- La espiral Misak.} El profe presenta el tiempo en espiral de los Misak (mirar atrás para avanzar) y pregunta: \emph{¿por qué mejorar mirando lo hecho lleva más lejos que empezar de cero cada vez?}. \textbf{Momento B (15 min) --- Autopsia de mi prototipo.} Cada quien saca su prototipo del módulo 3 y su tabla de 5 escenarios, y hace su \textbf{autopsia}: ¿cuáles pasó y cuáles no? Para cada falla, escribe el \textbf{síntoma} (\emph{qué hizo mal}) y \textbf{una causa probable} (\emph{umbral muy bajo, regla incompleta, sensor contaminado}). \textbf{Momento C (15 min) --- ¿A quién sirve, a quién no?} En parejas, cada uno le cuenta al otro para quién hizo su aparato, y juntos buscan \textbf{a quién NO le serviría} o en qué caso \textbf{fallaría de forma peligrosa}. \textbf{En el cuaderno:} tus fallas con síntoma y causa probable + una nota inicial de a quién sirve y a quién no.
\end{softbox}

{\sffamily\bfseries\fontsize{8}{10}\selectfont\textcolor{milcNegro!55}{MARCA CUANDO LO HAYAS HECHO}}
\milccheckrow{Puedo explicar por qué mejorar mirando lo hecho supera empezar de cero.}
\milccheckrow{Hice la autopsia de mi prototipo, con cada falla en su síntoma y causa probable.}
\milccheckrow{Escribí a quién le sirve mi aparato y a quién no (o dónde fallaría).}

\begin{infoband}{milcVerde}


\textbf{Lo que va al cuaderno (Actividad 1 · IDENTIFICA).} Título: «Actividad 1 --- Mirar atrás las fallas». \textbf{Formato:} tabla de fallas (síntoma → causa probable) de mis 5 escenarios + una nota de a quién sirve y a quién no mi aparato. \textbf{Extensión:} 8 líneas. \textbf{Sabes que terminaste cuando} tienes, para al menos una falla, \textbf{una causa probable concreta} que podrías poner a prueba.
\end{infoband}

\puente{Ya viste tus fallas. Ahora aprende el oficio de mejorar: depurar buscando la causa, el ciclo de una mejora por vuelta, evaluar con evidencia y reconocer límites.}

\milcestacion{milcTurquesa}{Estación 2 de 3 · Entender}

\begin{softbox}{milcTurquesa}{milcGris}{Lo que vas a entender}
\textbf{Actividad 2 · EXPLICA --- El oficio de mejorar} (55 min · grupo + individual). Ya viste tus fallas. Ahora aprende a convertirlas en una mejor versión: \textbf{depurar} buscando la causa, girar el \textbf{ciclo de mejora}, evaluar con \textbf{evidencia}, y reconocer \textbf{límites y ética}.
\end{softbox}

\milcpaso{1}{milcTurquesa}{\textbf{Pilar 1 --- Depurar es buscar la causa, no adivinar.} Cuando algo falla, el principiante cambia cosas al azar hasta que «funcione» ---y nunca sabe por qué. El investigador \textbf{sigue el agua hasta la gotera}: lee el \textbf{síntoma} (\emph{en el escenario 3 mostró sol estando oscuro}), plantea \textbf{una hipótesis} de causa (\emph{el umbral está muy bajo}), la \textbf{pone a prueba} (sube el umbral y vuelve al escenario 3), y si no era, plantea otra. \textbf{Una causa a la vez.} Depurar así no solo arregla el aparato: te \emph{enseña} cómo funciona. La falla, mirada de frente, es la mejor maestra.}
\milcpaso{2}{milcTurquesa}{\textbf{Pilar 2 --- El ciclo de mejora: una cosa por vuelta.} Mejorar es girar un lazo: \textbf{probar → observar dónde falla → cambiar UNA cosa → volver a probar}. La regla de oro está en el centro: \textbf{una sola cosa por vuelta}. Si cambias el umbral \emph{y} la regla \emph{y} el ícono a la vez y mejora, no sabrás cuál lo arregló ---ni podrás repetirlo. Cambia una, comprueba, anota; luego la siguiente. Es la espiral Misak hecha método: cada vuelta vuelve al aparato con \textbf{una} cosa aprendida, y sube un poco. \emph{Rehacer todo es tentador y casi siempre peor: pierdes lo que ya servía.}}
\milcpaso{3}{milcTurquesa}{\textbf{Pilar 3 --- Evaluar con evidencia: mide antes y después.} «Quedó mejor» no es una evaluación: es una opinión. Evaluar es \textbf{poner un número}. Elige un \textbf{criterio} claro ---\emph{¿cuántos de los 5 escenarios pasa? ¿en cuántos de 10 intentos acierta?}--- y mídelo \textbf{antes} de tu mejora y \textbf{después}. \emph{``Pasaba 3 de 5; cambié el umbral; ahora pasa 5 de 5''}: eso es una mejora \textbf{demostrada}, no sentida. Y si tras el cambio empeoró en otro escenario, la evidencia también te lo dirá ---y eso es oro, porque las mejoras a veces rompen lo que ya andaba. \textbf{Medir antes y después es lo que separa mejorar de creer que mejoraste.}}
\milcpaso{4}{milcTurquesa}{\textbf{Pilar 4 --- Límites y ética: qué no hace y a quién sirve.} Ningún aparato lo hace todo, y decirlo es de investigador maduro. Reconoce sus \textbf{límites}: qué condiciones no maneja (\emph{con poca luz se confunde}), qué tan probado está (\emph{solo 5 escenarios, un lugar}). Y su \textbf{ética}: \textbf{a quién sirve} (la abuela que oirá el timbre) y \textbf{a quién no} ---o dónde fallar sería \textbf{peligroso}. Una alarma \emph{ayuda} a vigilar, pero \textbf{no reemplaza} el cuidado humano; presentarla como una garantía sería engañar a quien confía en ella. \emph{Un aparato honesto dice, junto a lo que hace, lo que no hace ---sobre todo cuando alguien podría confiarle algo importante.}}

\begin{drawbox}{milcTurquesa}{Depurar + ciclo de mejora + evidencia + límites}
\textbf{Depurar:} síntoma → una hipótesis de causa → probarla → si no, otra. Una causa a la vez. \\[1mm]
\textbf{Ciclo de mejora:} probar → observar → cambiar UNA cosa → volver a probar. Una sola cosa por vuelta. \\[1mm]
\textbf{Evidencia:} elige un criterio y mídelo antes y después. «Pasaba 3 de 5, ahora 5 de 5», no «quedó mejor». \\[1mm]
\textbf{Límites y ética:} qué no hace · a quién sirve y a quién no · dónde fallar sería peligroso.
\end{drawbox}

\begin{softbox}{milcMostaza}{milcGris}{Errores comunes y su corrección}
\textbf{Errores al evaluar y mejorar (y cómo evitarlos).} (1) \emph{Cambiar varias cosas a la vez} --- si mejora, no sabes cuál fue; \textbf{una por vuelta}. (2) \emph{«Quedó mejor» sin medir} --- eso es opinión; \textbf{pon un número antes y después}. (3) \emph{Rehacer todo al primer fallo} --- pierdes lo que ya servía; \textbf{itera, no reinicies}. (4) \emph{Ocultar los límites} --- un aparato sin límites declarados es sospechoso; \textbf{dilos a la vista}. (5) \emph{Vender la alarma como garantía} --- prometer más de lo que hace puede poner a alguien en riesgo; \textbf{di a quién NO le basta}. (6) \emph{Adivinar en vez de depurar} --- cambiar al azar arregla sin enseñar; \textbf{busca la causa, una a la vez}.
\end{softbox}

\begin{infoband}{milcTurquesa}
\guiaDiagrama{assets/robotica-4/ciclo-de-mejora.tex}{El ciclo de mejora, la espiral Misak hecha método: cada vuelta cambia una sola cosa y la comprueba, para saber con evidencia qué la mejoró — no con suerte.}

\textbf{Lo que va al cuaderno (Actividad 2 · EXPLICA).} Título: «Actividad 2 --- El oficio de mejorar». \textbf{Formato:} 3 bloques. Bloque A: el ciclo de mejora dibujado (probar → observar → cambiar una cosa → volver a probar). Bloque B: mi criterio de evaluación y cómo lo mediré antes/después. Bloque C: la nota de límites y ética (qué no hace, a quién sirve, a quién no). \textbf{Extensión:} 1 hoja. \textbf{Sabes que terminaste cuando} podrías explicar \textbf{por qué se cambia una sola cosa por vuelta}.
\end{infoband}

\puente{Ya tienes las herramientas. Es hora de dar la vuelta de la espiral: iterar una mejora real, medir que sirvió, y escribir con honestidad los límites de tu aparato.}

\milcestacion{milcMagenta}{Estación 3 de 3 · Hacer}

\begin{softbox}{milcMagenta}{milcGris}{Cómo vas a construir tu producto}
\textbf{Actividad 3 · EVALÚA + CREA --- Da la vuelta de la espiral} (75 min · individual o parejas). Hora de mejorar de verdad. Vas a \textbf{depurar} una falla, \textbf{iterar} una mejora midiendo antes y después, y escribir la \textbf{nota honesta de límites} de tu aparato ---para cerrar la línea con un prototipo que evolucionó.
\end{softbox}

\milcpaso{1}{milcMagenta}{\textbf{Paso 1 --- Elige una falla y su causa (10 min).} De tu autopsia, elige \textbf{la falla que más importe} (la que rompe un escenario clave o la más peligrosa). Escribe su \textbf{síntoma} y \textbf{una hipótesis} de causa, concreta y comprobable: \emph{``en penumbra muestra sol; creo que el umbral está muy bajo''}. No dos hipótesis: una.}
\milcpaso{2}{milcMagenta}{\textbf{Paso 2 --- Mide el «antes» (10 min).} Antes de tocar nada, define tu \textbf{criterio} y mide el estado actual: \emph{``de 5 escenarios paso 3''}, o \emph{``de 10 intentos acierta 6''}. Anótalo. \textbf{Sin este número, no podrás demostrar que mejoraste} ---solo creerlo.}
\milcpaso{3}{milcMagenta}{\textbf{Paso 3 --- Cambia UNA cosa y vuelve a probar (20 min).} Aplica \textbf{un solo} cambio guiado por tu hipótesis (sube el umbral, completa la regla). Corre otra vez los escenarios y mide el \textbf{«después»}: \emph{``ahora paso 5 de 5''}. \textbf{¿Mejoró? ¿empeoró algo más?} Anota el resultado. Si no mejoró, tu hipótesis no era ---vuelve al Paso 1 con otra causa. Así gira la espiral.}
\milcpaso{4}{milcMagenta}{\textbf{Paso 4 --- Documenta la iteración (10 min).} Escribe tu mejora como \textbf{historia con evidencia}: \emph{``Falla: en penumbra mostraba sol. Causa: umbral en 20, muy bajo. Cambio: subí el umbral a 40. Antes: 3/5. Después: 5/5''}. Esa frase ---con números antes y después--- es la joya del módulo: prueba que mejoraste \textbf{con evidencia}, no con suerte.}
\milcpaso{5}{milcMagenta}{\textbf{Paso 5 --- Escribe la nota de límites y ética (15 min).} En un recuadro aparte: \textbf{Qué NO hace} mi aparato (condiciones que no maneja, qué tan probado está). \textbf{A quién le sirve} (la persona concreta) y \textbf{a quién NO le basta} (dónde haría falta algo más). \textbf{Dónde fallar sería grave} (si aplica) y por eso \textbf{no debe usarse como garantía}. \emph{Escribir esto no debilita tu trabajo: lo vuelve confiable y honesto, digno de la palabra caminada.}}
\milcpaso{6}{milcMagenta}{\textbf{Paso 6 --- Cierra la línea de Robótica (10 min).} Cierra tu cuaderno con la \textbf{evaluación liberadora} de toda la línea: \emph{¿Qué aprendí escuchando, diseñando, construyendo y mejorando? ¿Qué error me enseñó más? ¿A quién le sirve lo que hice? ¿Qué nueva pregunta me llevo?}. \emph{Esa pregunta nueva es el centro de tu próxima espiral} ---el arranque de otro itinerario, en robótica o en cualquier línea del semillero.}

\begin{drawbox}{milcMagenta}{Antes de cerrar tu prototipo evolucionado}
\checkbox Elegí una falla con una hipótesis de causa concreta y comprobable. \\[1mm]
\checkbox Medí el «antes» con un criterio claro (ej. paso 3 de 5). \\[1mm]
\checkbox Cambié una sola cosa, volví a probar y medí el «después». \\[1mm]
\checkbox Documenté la iteración con números antes/después. \\[1mm]
\checkbox Escribí la nota de límites y ética (qué no hace, a quién sirve y a quién no). \\[1mm]
\checkbox Cerré con mi evaluación liberadora de toda la línea y una pregunta nueva.
\end{drawbox}

\begin{softbox}{milcMagenta}{milcGris}{Plantilla de trabajo}
\textbf{Plantilla de la mejora.} \textit{Falla:} \_\_\_\_ · \textit{Causa (hipótesis):} \_\_\_\_. \\[1mm]
\textit{Criterio:} \_\_\_\_ · \textit{Antes:} \_\_/\_\_ · \textit{Cambio (una cosa):} \_\_\_\_ · \textit{Después:} \_\_/\_\_. \\[1mm]
\textit{Límites:} NO hace \_\_\_\_ · \textit{Sirve a} \_\_\_\_ · \textit{No le basta a} \_\_\_\_ · \textit{No usar como garantía si} \_\_\_\_.
\end{softbox}

\begin{infoband}{milcMagenta}


\textbf{Lo que va al cuaderno (Actividad 3 · EVALÚA + CREA).} Título: «Actividad 3 --- Di la vuelta de la espiral». \textbf{Formato:} la iteración documentada (falla → causa → antes → cambio → después, con números) + la nota de límites y ética + la evaluación liberadora de la línea con una pregunta nueva. \textbf{Extensión:} 1-2 hojas. \textbf{Sabes que terminaste cuando} tu mejora está \textbf{demostrada con números} y tu aparato dice, a la vista, lo que no hace.
\end{infoband}

\puente{El producto no es un aparato «perfecto»: es uno \textbf{evaluado y mejorado con evidencia}, con sus límites a la vista. Eso vale más que uno que nunca falló porque nunca se probó.}

\milcestacion{milcMostaza}{Producto de la sesión}
\textbf{Prototipo evolucionado — iteración con evidencia + nota de límites y ética}.
\begin{softbox}{milcMostaza}{milcGris}{Criterios de calidad}
(1) La depuración parte de un síntoma y una hipótesis de causa comprobable. (2) Se cambió una sola cosa por iteración (no varias a la vez). (3) La mejora está demostrada con un criterio medido antes y después (números). (4) Hay una nota de límites que declara qué no hace el aparato y a quién sirve y a quién no. (5) Cierra con una evaluación liberadora de la línea y una pregunta nueva.
\end{softbox}

\puente{Antes de cerrar, mírate como investigador de toda la línea: ¿qué aprendiste construyendo?, ¿qué error te enseñó más?, ¿a quién le sirve lo que hiciste?}

\milcestacion{milcVino}{Evaluación liberadora · cinco dimensiones}

\begin{softbox}{milcVino}{milcGris}{Sentido de la evaluación}
Evaluar no es solo revisar una nota. Es reconocer qué aprendí, qué cambió en mí, qué emociones manejé, cómo me relaciono con la ciudad y la comunidad, y qué le dejaré a quien viene detrás.
\end{softbox}

\textbf{Mi reflexión liberadora en cinco dimensiones.} Completa cada frase iniciada:

\small
\begin{tabularx}{\textwidth}{>{\bfseries\RaggedRight\arraybackslash}p{3.4cm}Y>{\RaggedRight\arraybackslash}p{4.6cm}}
\rowcolor{milcVino}\color{white}Dimensión & \color{white}\bfseries Mi respuesta (frame starter) & \color{white}\bfseries Algo que descubrí\\
1. Personal & Lo que más cambió en mí fue \linead{6.5cm} & \linead{4.2cm}\\
2. Emocional & La emoción más fuerte fue \linead{2.5cm} y la regulé \linead{3cm} & \linead{4.2cm}\\
3. Ciudadana & En democracia esto importa porque \linead{6cm} & \linead{4.2cm}\\
4. Local (Cartago) & En mi barrio sería útil para \linead{6.5cm} & \linead{4.2cm}\\
5. Intergeneracional & A mi abuela/abuelo le diría \linead{6.5cm} & \linead{4.2cm}\\
\end{tabularx}
\normalsize

\begin{infoband}{milcVino}
\textbf{Para la próxima sustentación.} Vuelve a esta tabla en seis meses. Verás que las respuestas cambian — no porque lo aprendido cambie, sino porque tú cambias. La evaluación liberadora es una conversación contigo a través del tiempo.
\end{infoband}


\puente{Cerramos con 3 voces que te ayudan a entender por qué corregirse, reconocer los límites y cuidar a quien usa tu aparato es la parte más madura de investigar.}

\milcestacion{milcGradeDark}{Triángulo de pensamiento · cierre filosófico}

\begin{softbox}{milcVino}{milcGris}{Dussel · lente del nosotros}
\textit{«Analéctico quiere indicar el hecho real humano por el que todo hombre, todo grupo o pueblo se sitúa siempre más allá (aná-) del horizonte de la totalidad.»}

{\footnotesize\itshape\color{milcNegro!70}--- Enrique Dussel · Filosofía de la liberación (1977)}

Dussel dice que siempre hay realidad \emph{más allá} de mi horizonte. Tu prototipo es una totalidad pequeña: probaste 5 escenarios, en un lugar, unos días. \textbf{Reconocer sus límites es admitir que hay muchísimos casos afuera de tu horizonte} ---los que no probaste, las personas a quienes no les basta. Por eso decir «a quién NO le sirve» no es rebajar tu aparato: es honrar lo que queda por fuera. \textbf{Un investigador honesto no confunde su totalidad con todo el mundo.}

\textbf{¿Reconozco lo que mi aparato deja afuera, o hablo como si mis 5 pruebas fueran todo el mundo?}
\end{softbox}
\linead{\linewidth}\\[1mm]
\linead{\linewidth}

\begin{softbox}{milcMostaza}{milcGris}{Estoicismo · Marco Aurelio · Meditaciones VI, 21 (c. 175 d.C.) · cuidado interior}
\textit{«Si alguien puede refutarme y probar de modo concluyente que pienso o actúo incorrectamente, de buen grado cambiaré de proceder. Pues persigo la verdad, que no dañó nunca a nadie; en cambio, sí se daña el que persiste en su propio engaño e ignorancia.»}

{\footnotesize\itshape\color{milcNegro!70}--- Marco Aurelio · Meditaciones VI, 21 (c. 175 d.C.)}

Marco Aurelio pone la verdad por encima de tener razón. Toda esta fase es eso: cuando un escenario \textbf{te refuta} ---muestra que tu aparato falla---, la respuesta madura no es defenderlo, es \textbf{cambiar de buen grado}. Persistir en un prototipo que falla «porque me costó hacerlo» es aferrarse al propio engaño. \textbf{Iterar es la forma técnica de amar la verdad más que el orgullo}: dejas que la prueba te corrija, y por eso tu aparato mejora.

\textbf{Cuando una prueba refuta mi aparato, ¿lo corrijo de buen grado, o lo defiendo por orgullo?}
\end{softbox}
\linead{\linewidth}\\[1mm]
\linead{\linewidth}

\begin{softbox}{milcTurquesa}{milcGris}{Floridi · lente de la infoesfera}
\textit{«El florecimiento de las entidades informacionales y de toda la infoesfera debe promoverse preservando, cultivando y enriqueciendo sus propiedades.»}

{\footnotesize\itshape\color{milcNegro!70}--- Luciano Floridi · La ética de la información (2013; trad.)}

Floridi pide cuidar el entorno común de la información, como se cuida un río. Tu aparato entra al mundo de otras personas: si prometes que una alarma «garantiza» algo que no garantiza, \textbf{contaminas} la confianza de quien la usa ---y el daño puede ser real. Decir con honestidad a quién sirve y a quién no, qué hace y qué no, \textbf{enriquece} ese entorno común en vez de ensuciarlo. \textbf{Cuidar a quien usará tu aparato es parte de construirlo bien.}

\textbf{¿Mi aparato, y lo que digo de él, cuidan a quien lo va a usar, o le prometen de más?}
\end{softbox}
\linead{\linewidth}\\[1mm]
\linead{\linewidth}

\begin{infoband}{milcNegro}
\textbf{Nota para el docente.} Las tres son citas textuales y llevan su referencia debajo. Puedes remitir a la obra en clase.
\end{infoband}

\milcestacion{milcVino}{Mi autoevaluación · marca una casilla por fila}

{\small
\setlength{\extrarowheight}{2.2pt}
\begin{tabularx}{\textwidth}{>{\RaggedRight\arraybackslash\bfseries}p{3.0cm}YYY}
\rowcolor{milcVino}\color{white}\bfseries Criterio & \color{white}\bfseries Lo logré & \color{white}\bfseries En proceso & \color{white}\bfseries Necesito apoyo\\[.6mm]
Comprensión del tema & \milcopcion{Explico las ideas con mis palabras.} & \milcopcion{Reconozco las ideas, falta precisión.} & \milcopcion{Necesito repasar lo básico.}\\[.8mm]
\rowcolor{milcGris}Pensamiento computacional & \milcopcion{Usé los pilares al armar mi producto.} & \milcopcion{Usé algunos pilares.} & \milcopcion{Necesito releer la estación 2.}\\[.8mm]
Aplicación y producto & \milcopcion{Mi producto cumple los criterios.} & \milcopcion{Está armado, con ajustes pendientes.} & \milcopcion{Aún está en construcción.}\\[.8mm]
\rowcolor{milcGris}Reflexión 5 dimensiones & \milcopcion{Respondí cada una con honestidad.} & \milcopcion{Respondí algunas con detalle.} & \milcopcion{Mis respuestas son muy generales.}\\[.8mm]
Triángulo de pensamiento & \milcopcion{Conecté las 3 voces con mi trabajo.} & \milcopcion{Conecté al menos una.} & \milcopcion{No las relaciono con mi proceso.}\\[.8mm]
\end{tabularx}}

\vspace{3mm}
\textbf{Hoy entendí que:}\\[1mm]
\linead{\linewidth}\\[1mm]
\linead{\linewidth}

\textbf{Mi compromiso para la próxima sesión es:}\\[1mm]
\linead{\linewidth}\\[1mm]
\linead{\linewidth}

\begin{softbox}{milcMostaza}{milcGris}{Cita de despedida · Marco Aurelio}
\textit{``El obstáculo en el camino se vuelve el camino. Lo que estorba al camino se vuelve camino.''}\\[1mm]
Cada error de hoy, bien observado, se convierte en pieza del camino que pisarás mañana.
\end{softbox}

\small
\begin{tabularx}{\textwidth}{>{\bfseries}p{3.6cm}Y}
\rowcolor{milcGradeDark!12}Nombre completo: & \linead{10cm}\\
Fecha: & \linead{4cm} \quad Curso/grupo: \linead{4cm}\\
Mi firma: & \linead{9cm}\\
\end{tabularx}
\normalsize

\begin{infoband}{milcGradeDark}
\textbf{Para cerrar tu proyecto de Robótica.} Dos cosas. \textbf{(1) Deja tu prototipo evolucionado y su historia:} el registro de fallas y depuración, la iteración con números antes/después, la versión mejorada, la nota de límites y ética, y el video. Si aún quedan fallas, no las escondas: anótalas como \textbf{trabajo pendiente} ---la espiral sigue girando. \textbf{(2) Devuélvele el aparato a su dueño:} entrégaselo (o muéstraselo) a la persona a quien le sirve, cuéntale con honestidad qué hace y qué no, y déjalo probarlo. \textbf{Trae:} tu proyecto completo de la línea ---escucha, diseño, construcción, mejora--- listo para exponer en el aula o en una feria. \emph{Con esto completas tu segundo itinerario del semillero, de principio a fin. Miraste el reclamo de una necesidad, diseñaste su lógica, construiste el aparato y lo hiciste mejor con evidencia. Como los Misak, no terminaste en línea recta: cerraste una vuelta de la espiral. La pregunta nueva que anotaste es el centro de la siguiente ---en robótica, o en la línea que tu curiosidad elija.}
\end{infoband}

\end{document}
